\scanpage{124}
\mspage{\handwritten{2}}
% Three long diagonal pen strokes cross the typed page.
\deleted{$Z(S')={}$ ensemble des couples $(g,v)\handwritten{\in}G(S')\times V(S')$ tels que $v\handwritten{\in}g.W_{S'}(S')$. Par suite, $Z$ est isomorphe au faisceau $G\times_S W$, grâce à l'isomorphisme
\[
G\times_S W\handwritten{\xrightarrow{\sim}}Z
\]
défini par $(g,w)\handwritten{\rightsquigarrow}(g,g.w)$. Ainsi le diagramme cartésien précédent donne le diagramme cartésien
\diagram{\begin{tikzcd}[ampersand replacement=\&,column sep=large]
\& G\times_S W\handwritten{\arrow[r,"q"]\arrow[d,"\lambda"']}\&X\handwritten{\arrow[d,"i"]}\\
\handwritten{(\mathrm{xx})}\&G\times_S V\handwritten{\arrow[r]}\&G/N\times_S V
\end{tikzcd}}
où on a $\handwritten{\lambda}(g,w)=(g,g.w)$, donc $q(g,w)=(g^{\bullet},g.w)$, où $g^{\bullet}$ désigne l'image de $g$ par l'application canonique $G(S')\handwritten{\to}(G/N)(S')$. On voit enfin sur le diagramme (x) que $Z\handwritten{\to}X$ fait de $Z$ un fibré principal \deleted{\unsure{homogène} sur} \inserted{de base} $X$ \inserted{et} de groupe $N$ opérant \deleted{\unsure{par}} à droite par $(g,v).n=(gn,v)$, de sorte que \deleted{dans} dans (xx), $q\colon G\times_S W\handwritten{\to}X$ fait de $G\times_S W$ un fibré principal de base $X$ et de groupe $N$ opérant \deleted{\unsure{par}} à droite par
\[
(g,w).n=(gn,n^{-1}.w)\,.
\]
No\unsure{n|u}s \deleted{\unsure{pouvons}} \handwritten{résumons} \deleted{considérer} \inserted{\handwritten{les principaux morphismes précédents dans}} le diagramme suivant \handwritten{:}
\diagram{\begin{tikzcd}[ampersand replacement=\&,column sep=large]
\&\&G\times_S W\handwritten{\arrow[d,"q"]\arrow[dr,"\varphi"']}\&\\
\handwritten{(\mathrm D)}\&G/N\&X\handwritten{\arrow[l,"p"]\arrow[r,"\psi"']\arrow[d,"i"]}\&V\\
\&\&G/N\times_S V\handwritten{\arrow[ul,dashed,"\operatorname{pr}_1"]\arrow[ur,dashed,"\operatorname{pr}_2"']}\&
\end{tikzcd}}
où $\handwritten{\psi}=\operatorname{pr}_2\,i$ et $\unsure{\phi|\varnothing}=\handwritten{\psi}\,q$ \inserted{i.e. $\phi(g,w)=g.w$.} Si $\deleted{\unsure{u}}\handwritten{v}$ est une section de $V$ sur $S$, le sous-faisceau $\deleted{\unsure{P|F}}\handwritten{X}_v$ de $X$ image inverse de cette section par $\handwritten{\psi}$ est donné par

$\deleted{\unsure{A|P}}\handwritten{X}_v(S')={}$ ensemble des sous-faisceaux $W'$ de $V_{S'}$ qui sont localement conjugués de $W_{S'}$ par $G$, et qui contiennent la section $v_{S'}$ de $V_{S'}$, tandis que le sous-faisceau $\deleted{\unsure{X_v|Z_v}}$ de $G\times_S W$ image inverse de $v$ par $\phi$ est \deleted{donné par} isomorphe au sous-faisceau $M_v$ de $G$ \deleted{\unsure{formé}} donné par}
% slanted
\margin{\handwritten{\unsure{lu|$M_v$}}}
% archive sheet number: 123
